\documentclass[12pt]{article}


\usepackage[margin=2cm]{geometry}
\usepackage[shortlabels]{enumitem}
\usepackage{amsmath,amsxtra,amsthm,amssymb,mathrsfs}
\usepackage[all]{xy}

\usepackage{exsheets}
\SetupExSheets{headings=runin}
\SetupExSheets{question/type=exercise}
\SetupExSheets{question/name=Ex }

\newcommand{\Hom}{\operatorname{Hom}\nolimits}
\newcommand{\End}{\operatorname{End}\nolimits}

\DeclareMathOperator{\Char}{char}
\newcommand{\triv}{\operatorname{triv}\nolimits}
\newcommand{\sgn}{\operatorname{sgn}\nolimits}
\newcommand{\Mat}{\operatorname{Mat}\nolimits}
\newcommand{\GL}{\operatorname{GL}\nolimits}
\DeclareMathOperator{\Sn}{\mathfrak{S}}

\begin{document}

{\scshape Autumn 2025} \hfill {\scshape \large Topics in Math. Sci. VIII} \hfill {\scshape Homework Assignment 1}
\smallskip

\hrule

\bigskip

\begin{question}
Let $A=A_1\times A_2$ be a direct product of two $\Bbbk$-algebras $A_1, A_2$.
Note that there are algebra homomorphisms $\pi_i : A\to A_i$.
Consider the idempotents $e_1:=(1_{A_1},0), e_2:=(0,1_{A_2})$ of $A$.
\begin{enumerate}
\item For $M\in \mod A$, show that $M=M_1\oplus M_2$ with $M_i=Me_i$ for both $i\in\{1,2\}$, i.e. $M=Me_1+Me_2$ with $Me_1\cap Me_2=0$.

\item Show that $\pi_i$ induces a simple $A$-module structure on the simple $A_i$-modules for both $i\in\{1,2\}$.

\item Show that there exists a natural $A_i$-action on the direct summand $M_i$ of $M$ (in the notation of (1)).  In particular, classify the simple $A$-modules.
\end{enumerate}
\end{question}

%\begin{question}
%Consider the formal power series ring $\Bbbk[[x]]:=\{\sum_{i=0}^\infty \lambda_i x^i\mid \lambda_i\in \Bbbk\}$ with 
%\begin{itemize}
%\item addtition $(\sum_i \lambda_i x^i)+(\sum_i \mu_i x^i) = \sum_i (\lambda_i+\mu_i)x^i$,
%\item scalar multiplication $\lambda(\sum_i\lambda_ix^i)=\sum_i (\lambda \lambda_i)x^i$,
%\item multiplication $(\sum_i \lambda_i x^i)(\sum_j\mu_j x^j ) = \sum_k( \sum_{i+j=k} \lambda_i\mu_j) x^k$.
%\end{itemize}
%\begin{enumerate}
%\item Show that $\Bbbk[[x]]$ is a commutative $\Bbbk$-algebra.
%
%\item Determine the invertible elements in $\Bbbk[[x]]$.
%
%\item Show that $\Bbbk[[x]]$ is a local algebra, i.e. there exists a unique maximal left (equivalently, right) ideal.
%
%\item Classify the simple modules of $\Bbbk[[x]]$.
%
%\item Classify the simple modules of $\Bbbk[x]$.  \emph{Hint}: (a) Maximal two-sided ideals of a ring are determined by the irreducible elements. (b) $\Bbbk[x]$ is commutative.
%\end{enumerate}
%\end{question}
%
%\begin{question}
%Consider the group algebra $A=R[G]$ of the finite group $G=C_2=\{1,\sigma\}$.
%\begin{enumerate}
%\item For $R=\mathbb{C}$, show that there are two non-isomorphic simple $A$-modules $S_1,S_2$ so that $A\cong S_1\oplus S_2$.  Hint: Find an a non-zero non-identity idempotent of $A$.
%
%\item For $R=\mathbb{F}_2$ the finite field with two elements, show that $A$ is a local algebra by, for example, showing that $A\cong R[x]/(x^2)$ as $R$-algebra.
%\end{enumerate}
%\end{question}

\begin{question}
Let $A$ be the ring
\[
\begin{pmatrix}
\mathbb{R} & \mathbb{C} \\ 0 & \mathbb{C}
\end{pmatrix} = \left\{ \begin{pmatrix} a&b\\0&c \end{pmatrix} : a\in \mathbb{R}, b,c\in \mathbb{C} \right\}.
\]
\begin{enumerate}
\item Show that $A$ is an $\mathbb{R}$-algebra and find the dimension of $A$ over $\mathbb{R}$.

\item There are two simple $A$-modules.  Describe them.

\item Write down a composition series (i.e. filtration whose subquotients are simples) of the free $A$-module $A_A$.
\end{enumerate}
\end{question}


\begin{question}
Let $A$ be the following $\Bbbk$-algebra:
\[
\left\{
\begin{pmatrix}
a & b & c \\ 0& x & y \\ 0 & 0& a
\end{pmatrix} \mid a,b,c,x,y\in \Bbbk
\right\}
\]
Note that for every matrix in $A$, the (1,1)-entry and the (3,3)-entry must be the same.
Let $e_1$ be the idempotent of $A$ given by the matrix with 1 in the (1,1)- and (3,3)-entry and 0 everywhere else; and $e_2=1_B-e_1$.

\begin{enumerate}
\item Show that both $e_1A$ and $e_2A$ have a unique simple submodules and they are isomorphic.

\item Let $S_1$ be the simple module in (1).  Show that $S_2:=e_2A/S_1 \ncong S_1$.

\item Find the composition series of $e_1A$ and of $e_2A$.
\end{enumerate}
\end{question}


\begin{question}
\begin{enumerate}
\item Consider the linearly oriented $\vec{\mathbb{A}}_5$-quiver:
\[
\xymatrix{ 1 \ar[r]^{\alpha_1} & 2 \ar[r]^{\alpha_2} & 3 \ar[r]^{\alpha_3} & 4 \ar[r]^{\alpha_4} & 5}
\]
and the following representation $M$ of $\vec{\mathbb{A}}_5$:
\[
\Bbbk \xrightarrow{1} \Bbbk \xrightarrow{\left(\begin{smallmatrix}1\\0\end{smallmatrix}\right)} \Bbbk^2 \xrightarrow{\left(\begin{smallmatrix}a&b\\c&d\end{smallmatrix}\right)} \Bbbk^2 \xrightarrow{(1,1)} \Bbbk 
\]
Find the indecomposable decompositions of $M$ in the cases when the matrix $\left(\begin{smallmatrix}a&b\\c&d\end{smallmatrix}\right)$ is of rank 1 and of rank 2.
Explain in details your reasoning.

You may use the fact that there are indecomposable modules of the form $U_{i,j}$ for $1\leq i\leq j\leq 5$ such that
\[
U_{i,j}e_x = \begin{cases}
\Bbbk & \text{if  $i\leq x\leq j$;}\\
0 & \text{otherwise,}
\end{cases} \text{ and } U_{i,j}\alpha_k = \begin{cases}
\mathrm{id} & \text{if $i\leq k<j$;}\\
0 & \text{otherwise.}
\end{cases}
\]

\item Consider $A=\Bbbk Q/I$ given by
\[
Q: \xymatrix@1{1 \ar[r]^{\alpha}\ar@/_1pc/[rr]_{\delta}& 2 \ar[r]^{\beta}& 3 \ar[r]^{\gamma}& 4}, \quad I=\langle \alpha\beta\gamma-\delta\gamma\rangle 
\]

\begin{enumerate}
\item Find a basis for the socle (i.e. maximal semisimple submodule) of the indecomposable projective $P_1$.\\
\emph{Hint}: The socle is 2-dimensional.

\item Show that the radical $\mathrm{rad}P_1:=\mathrm{Ker}(P_1\twoheadrightarrow S_1)$ of $P_1$ is not indecomposable.  Details your reason.
\end{enumerate}

%\item Consider the following quiver 
%\[
%Q:\quad  \xymatrix{ 1 \ar@(dl,ul)^{\alpha} \ar@/^/[r]^{\beta}& 2 \ar@/^/[l]^{\gamma}}
%\]
%Let $I_1:=\langle \alpha^2-\beta\gamma, \gamma\beta-\gamma\alpha\beta,\alpha^4\rangle$ and $I_2:=\langle \alpha^2-\beta\gamma,\gamma\beta,\alpha^4\rangle$.
%Show that $\Bbbk Q/I_1 \cong \Bbbk Q/I_2$ as $\Bbbk$-algebra when the characteristic of $\Bbbk$ is not 2.
\end{enumerate}
\end{question}



%\begin{question}
%Let $A$ be the algebra of upper triangular $n\times n$-matrices:
%\[
%A := \begin{pmatrix}
%K & K & \cdots & K \\
%0 & K & \cdots & K \\
%0 & 0 & \ddots & \vdots \\
%0 & 0 & 0 & K
%\end{pmatrix} = { \left\{(a_{i,j})_{1\leq i,j\leq n}\middle| \begin{array}{l}a_{i,j}\in K\;\forall i,j\\
%a_{i,j}=0\;\forall i>j \end{array}\right\} }
%\]
%For $1\leq i \leq j \leq n$, let $M_{i,j} \subset K^{\oplus n}$ be the vector space given by column vectors $v=(v_k)_{1\leq k\leq n}$ where $v_k=0$ for $k\notin \{i,i+1,\ldots, j\}$.
%\begin{enumerate}[\rm(i)]
%\item Determine which $M_{i,j}$'s are simple.
%\item Describe the composition series of $M_{i,j}$.
%\end{enumerate}  
%\end{question}

\vfill
\hfill Deadline: 24th October, 2025

\hfill Submission: E-mail to (replace at by @) aaron.chan at math.nagoya-u.ac.jp


\end{document}